\documentclass{tufte-handout}

%\geometry{showframe}% for debugging purposes -- displays the margins

\usepackage[spanish, es-tabla]{babel}
\usepackage{amsmath}

% Set up the images/graphics package
\usepackage{graphicx}
\setkeys{Gin}{width=\linewidth,totalheight=\textheight,keepaspectratio}
\graphicspath{{graphics/}}

\title[Química Física II: Ejercicios Tema 4]{
Química Física II: Ejercicios Tema 4}
%\author{David De Sancho}
\date{}  % if the \date{} command is left out, the current date will be used

% The following package makes prettier tables.  We're all about the bling!
\usepackage{booktabs}

% The units package provides nice, non-stacked fractions and better spacing
% for units.
\usepackage{units}

% The fancyvrb package lets us customize the formatting of verbatim
% environments.  We use a slightly smaller font.
\usepackage{fancyvrb}
\fvset{fontsize=\normalsize}

% Small sections of multiple columns
\usepackage{multicol}

% Provides paragraphs of dummy text
\usepackage{lipsum}

% These commands are used to pretty-print LaTeX commands
\newcommand{\doccmd}[1]{\texttt{\textbackslash#1}}% command name -- adds backslash automatically
\newcommand{\docopt}[1]{\ensuremath{\langle}\textrm{\textit{#1}}\ensuremath{\rangle}}% optional command argument
\newcommand{\docarg}[1]{\textrm{\textit{#1}}}% (required) command argument
\newenvironment{docspec}{\begin{quote}\noindent}{\end{quote}}% command specification environment
\newcommand{\docenv}[1]{\textsf{#1}}% environment name
\newcommand{\docpkg}[1]{\texttt{#1}}% package name
\newcommand{\doccls}[1]{\texttt{#1}}% document class name
\newcommand{\docclsopt}[1]{\texttt{#1}}% document class option name

\begin{document}
\maketitle
\large 

\begin{enumerate}
 \item Comprueba que el armónico esférico $Y_{10}$ es una solución 
  de la ecuación $\Lambda^2Y_{l,m_l}=-l(l+1)Y_{l,m_l}$.
% \item Una esfera de 250 g y 4 cm de radio gira a 5 revoluciones por 
% segundo. Estima el valor de $l$ y el ángulo mínimo correspondiente
% al vector momento angular.
\item Demuestra que el armónico esférico $Y_{10}$ es una función de onda normalizada.
  
\item Comprueba que el armónico esférico $Y_{1,-1}$ no es función propia del operador $\hat{l}_x$.
\end{enumerate}

\newpage
\subsection{\textbf{Soluciones:}}
\begin{enumerate}
    \item El valor propio del operador es -2, que en
  efecto es igual a $-l(l+1)$ para $l=1$.
  \item --
  \item --
\end{enumerate}

\end{document}